\subsection{Multiplicities and flat extensions}

PROPOSITION 4. — Let \(\rho: A \to B\) be a local homomorphism of Noetherian local rings, and let \(N\) be a \(B\) -module, flat over \(A\) , and such that \(N \otimes_A \kappa_A\) is a \(B\) -module of finite length. If \(M\) is a nonzero finitely generated \(A\) -module and let q be an ideal of \(A\) distinct from \(A\) and such that \(M/qM\) is of finite length, then \((M \otimes_A N)/(qB)(M \otimes_A N)\) is a nonzero finitely generated \(B\) -module of finite length, and we have

\[
e _ {\mathrm{qB}} ^ {\mathrm{B}} (\mathbf {M} \otimes_ {\mathrm{A}} \mathbf {N}) = \operatorname{long} _ {\mathrm{B}} (\mathbf {N} \otimes_ {\mathrm{A}} \kappa_ {\mathrm{A}}). e _ {\mathrm{q}} ^ {\mathrm{A}} (\mathbf {M}).
\]

Let L be an A-module of finite length r. Then L possesses a Jordan-Hölder sequence of length r, with quotients isomorphic to \(\kappa_{A}\) ; since N is flat over A, the B-module \(L \otimes_{A} N\) possesses a composition sequence of length r, with quotients isomorphic to \(N \otimes_{A} \kappa_{A}\) , hence is of length \(r \cdot \operatorname{long}_{\mathbf{B}}(\mathbf{N} \otimes_{\mathbf{A}} \kappa_{\mathbf{A}})\). Since the B-module \((\mathbf{M} \otimes_{\mathbf{A}} \mathbf{N})/(\mathbf{q}\mathbf{B})^{n}(\mathbf{M} \otimes_{\mathbf{A}} \mathbf{N})\) is isomorphic to \((\mathbf{M}/\mathbf{q}^{n}\mathbf{M}) \otimes_{\mathbf{A}} \mathbf{N}\) for every \(n \in \mathbf{N}\), the proposition follows from the definition of multiplicities.

COROLLARY. — Suppose that B is flat over A and that \(\rho(\mathfrak{m}_{\mathbf{A}})\) B = \(\mathfrak{m}_{\mathbf{B}}\) . Then

\[
e _ {\mathbf {q B}} ^ {\mathbf {B}} (\mathbf {M} \otimes_ {\mathbf {A}} \mathbf {B}) = e _ {\mathbf {q}} ^ {\mathbf {A}} (\mathbf {M}).
\]

This applies in particular when B is the completion * or the henselization * of A with respect to an ideal distinct from A, * or an inflation of A, for example a strict henselization of A. *

Example. — * Let X be a complex algebraic variety, \(\mathcal{O}_{X,x}\) the local ring of X at a rational point \(x\) , \(X^{an}\) the analytic space associated with X; denote again by \(x\) the point of \(X^{an}\) corresponding to \(x\) , and let \(\mathcal{O}_{X^{an},x}\) be the local ring of \(X^{an}\) at \(x\) . Then \(e(\mathcal{O}_{X^{an},x}) = e(\mathcal{O}_{X,x})\) . *

